\section{Eksisterende løsninger og valg av tilnærming}
\label{sec:eksisterende-loesninger}

Problemet VideoAPI løser kan angripes på flere måter. Dette delkapittelet
gjennomgår tre tilnærminger og forklarer hvorfor provider-basert
mellomvare ble vurdert som den mest hensiktsmessige løsningen for
HDOs behov.

Et beslektet eksempel finnes hos Singh~\cite{singh2026video}, som
foreslår en generisk WebRTC-komponent for å redusere koblingen mellom
webapplikasjoner og leverandørspesifikk signaleringslogikk. Forskjellen
er at denne løsningen plasserer abstraksjonen i klientlaget, mens
VideoAPI plasserer den i et server-side REST-mellomlag.

HDO benyttet direkte integrasjon som utgangspunkt, der hvert
klientsystem integrerer mot AntMedia og NHN separat~\cite{hdo_oppgave}.
Dette eliminerer et ekstra lag, men betyr at autentisering,
feilhåndtering og domeneoversettelse må implementeres separat per
motor, og at alle klienter må oppdateres dersom én motor byttes ut.
Det var nettopp denne dupliseringen og koblingen til konkrete
videomotorer som motiverte VideoAPI-prosjektet.

API-gateways som Kong og Azure API
Management~\cite{microsoft-api-gateway} sentraliserer sikkerhet og
observability, men er primært proxy- og policy-verktøy. De ruter
forespørsler uten å oversette mellom domenemodeller. Provider-basert
mellomvare med adapter-mønsteret~\cite{cockburn-hexagonal,
microsoft-clean-architecture} fungerer samtidig som et
anti-corruption layer~\cite{microsoft-acl} mellom HDOs kjernemodell
og de ulike videomotorenes semantikk.

\subsection*{Tilnærming 1: Direkte integrasjon}

\begin{figure}[H]
\centering
\begin{tikzpicture}[
    box/.style={
        rectangle, draw=black, thick, rounded corners=3pt,
        minimum width=3.2cm, minimum height=1.0cm,
        align=center, font=\small\ttfamily
    },
    motor/.style={box, fill=orange!15},
    client/.style={box, fill=blue!8},
    arr/.style={-{Stealth[length=5pt]}, thick},
    lbl/.style={font=\footnotesize\rmfamily\itshape}
]

\node[client] (k1) at (0,  0)   {Klient A};
\node[client] (k2) at (0, -2.5) {Klient B};
\node[motor]  (am) at (8,  0)   {AntMedia\\Server};
\node[motor]  (nhn) at (8, -2.5){NHN Video};

\draw[arr] (k1.east) -- (am.west)
    node[midway, above, lbl] {direkte kall};
\draw[arr] (k1.east) to[out=-20, in=160] (nhn.west)
    node[midway, below, lbl] {};
\draw[arr] (k2.east) to[out=20, in=200] (am.west)
    node[midway, above, lbl] {};
\draw[arr] (k2.east) -- (nhn.west)
    node[midway, below, lbl] {direkte kall};

\end{tikzpicture}
\caption{Direkte integrasjon: hver klient integrerer mot begge
videomotorene separat. Endringer i én motor påvirker alle klienter.}
\label{fig:tilnaerming-direkte}
\end{figure}

\subsection*{Tilnærming 2: API-gateway}

\begin{figure}[H]
\centering
\begin{tikzpicture}[
    box/.style={
        rectangle, draw=black, thick, rounded corners=3pt,
        minimum width=3.2cm, minimum height=1.0cm,
        align=center, font=\small\ttfamily
    },
    motor/.style={box, fill=orange!15},
    client/.style={box, fill=blue!8},
    gateway/.style={box, fill=gray!20, minimum width=2.8cm},
    arr/.style={-{Stealth[length=5pt]}, thick},
    lbl/.style={font=\footnotesize\rmfamily\itshape}
]

\node[client]  (k1)  at (0,  0)    {Klient A};
\node[client]  (k2)  at (0, -2.5)  {Klient B};
\node[gateway] (gw)  at (5, -1.25) {API\\Gateway};
\node[motor]   (am)  at (10,  0)   {AntMedia\\Server};
\node[motor]   (nhn) at (10, -2.5) {NHN Video};

\draw[arr] (k1.east) -- (gw.west);
\draw[arr] (k2.east) -- (gw.west);
\draw[arr] (gw.east) to[out=20,  in=180] (am.west)
    node[midway, above, lbl] {};
\draw[arr] (gw.east) to[out=-20, in=180] (nhn.west);

\node[font=\footnotesize\rmfamily, text=red!70!black]
    at (5, -3.2) {Ingen semantisk oversettelse};

\end{tikzpicture}
\caption{API-gateway: ett felles inngangspunkt og sentralisert
sikkerhet, men gatewayen ruter kun forespørsler videre.
Domenemismatch mellom AntMedia og NHN løses ikke.}
\label{fig:tilnaerming-gateway}
\end{figure}

\subsection*{Tilnærming 3: Provider-basert mellomvare}

\begin{figure}[H]
\centering
\resizebox{\linewidth}{!}{%
\begin{tikzpicture}[
    box/.style={
        rectangle, draw=black, thick, rounded corners=3pt,
        minimum width=3.2cm, minimum height=1.0cm,
        align=center, font=\small\ttfamily
    },
    motor/.style={box, fill=orange!15},
    client/.style={box, fill=blue!8},
    api/.style={box, fill=green!12, minimum width=2.8cm},
    provider/.style={box, fill=green!8},
    arr/.style={-{Stealth[length=5pt]}, thick},
    lbl/.style={font=\footnotesize\rmfamily\itshape}
]

\node[client]   (k1)   at (0, -1.25)  {Klient};
\node[api]      (api)  at (4.5,-1.25) {VideoAPI\\(felles API)};
\node[provider] (ant)  at (9,  0)     {AntMedia\\Provider};
\node[provider] (nhn)  at (9, -2.5)   {NHN\\Provider};
\node[motor]    (am)   at (13.5, 0)   {AntMedia\\Server};
\node[motor]    (nhnb) at (13.5,-2.5) {NHN Video};

\draw[arr] (k1.east)  -- (api.west);
\draw[arr] (api.east) to[out=20,  in=180] (ant.west);
\draw[arr] (api.east) to[out=-20, in=180] (nhn.west);
\draw[arr] (ant.east) -- (am.west);
\draw[arr] (nhn.east) -- (nhnb.west);

\node[font=\footnotesize\rmfamily, text=green!60!black]
    at (4.5,-3.2) {Semantisk oversettelse i provider-laget};

\end{tikzpicture}%
}
\caption{Provider-basert mellomvare: klienten forholder seg til ett
felles API. Provider-laget oversetter mellom felles modell og
motorspesifikk semantikk.}
\label{fig:tilnaerming-videapi}
\end{figure}

\begin{table}[H]
\centering
\caption[Sammenligning av tilnærminger]{Sammenligning av tilnærminger for integrasjon mot heterogene videomotorer.}
\label{tab:eksisterende-loesninger}
\setlength{\extrarowheight}{6pt}
\small
\begin{tabularx}{\textwidth}{>{\bfseries}X c c c}

\arrayrulecolor{black}
\toprule
\rowcolor{gray!25}
\textbf{Kriterium} &
\textbf{Direkte integrasjon} &
\textbf{API-gateway} &
\textbf{VideoAPI} \\

\rowcolor{gray!25}
& \textbf{\small(HDOs utgangspunkt)} &
\textbf{\small(Kong, Azure)} &
\textbf{\small(valgt løsning)} \\
\midrule

\rowcolor{white}
Felles API &
\textcolor{red!70!black}{\textbf{Nei}} &
\textcolor{green!60!black}{\textbf{Ja}} &
\textcolor{green!60!black}{\textbf{Ja}} \\

\rowcolor{gray!10}
Semantisk oversettelse &
\textcolor{red!70!black}{\textbf{Nei}} &
\textcolor{red!70!black}{\textbf{Nei}} &
\textcolor{green!60!black}{\textbf{Ja}} \\

\rowcolor{white}
Flermotor-støtte &
\textcolor{red!70!black}{\textbf{Nei}} &
\textcolor{orange!70!black}{\textbf{Delvis}} &
\textcolor{green!60!black}{\textbf{Ja}} \\

\rowcolor{gray!10}
Provider-utskiftbarhet &
\textcolor{red!70!black}{\textbf{Nei}} &
\textcolor{red!70!black}{\textbf{Nei}} &
\textcolor{green!60!black}{\textbf{Ja}} \\

\rowcolor{white}
Lav klientkobling &
\textcolor{red!70!black}{\textbf{Nei}} &
\textcolor{orange!70!black}{\textbf{Delvis}} &
\textcolor{green!60!black}{\textbf{Ja}} \\

\rowcolor{gray!10}
Selvhostet &
\textcolor{green!60!black}{\textbf{Ja}} &
\textcolor{orange!70!black}{\textbf{Delvis}} &
\textcolor{green!60!black}{\textbf{Ja}} \\

\rowcolor{white}
Lisenskostnad &
Ingen &
Varierer &
Ingen \\

\bottomrule
\end{tabularx}

\vspace{0.5em}
\begin{minipage}{\textwidth}
\footnotesize
\textit{Felles API}: klientene forholder seg til ett grensesnitt uavhengig av underliggende videomotor.\\
\textit{Semantisk oversettelse}: løsningen håndterer ulike domenemodeller mellom videomotorene eksplisitt.\\
\textit{Flermotor-støtte}: løsningen kan rute mot flere videomotorer samtidig.\\
\textit{Provider-utskiftbarhet}: en videomotor kan byttes eller legges til uten endringer i controller-laget.\\
\textit{Lav klientkobling}: klientene trenger ikke kjenne til motorspesifikke API-er, autentiseringsmekanismer eller datamodeller.\\
\textit{Selvhostet}: løsningen driftes i HDOs egen infrastruktur uten ekstern skytjeneste.\\
\textit{Lisenskostnad}: kostnad knyttet til bruk av løsningen i drift. For API-gatewayer varierer dette mellom åpne og kommersielle produkter.
\end{minipage}
\end{table}

Tabell~\ref{tab:eksisterende-loesninger} viser at direkte integrasjon
ikke gir ett felles API og krever duplisert implementasjon per motor.
API-gateways løser flere tverrgående infrastrukturbehov, men håndterer
ikke at AntMedias broadcast-semantikk og NHNs booking-semantikk er
fundamentalt ulike. Provider-basert mellomvare gir derfor en bedre
balanse i HDOs kontekst: klientene får ett felles API, mens
motorspesifikk autentisering, feilhåndtering og domenemapping holdes
isolert i provider-laget. % Hvorvidt arkitekturen faktisk oppfyller kravene fra HDO evalueres i delkapittel~\ref{sec:evaluering}.